%---------------------------------------------------------------------
% The fourteen questions.
%
% Each carries its provenance: (بہار 2013، سوال n) or (خزاں 2013، سوال n).
% Where a topic appeared in both papers, both references are given --
% that is the signal that it is high-yield.
%---------------------------------------------------------------------

\section*{صحافت کی بنیادیں}
\markboth{صحافت کی بنیادیں}{}

\begin{sawal}{سوال \textenglish{1} \ \textnormal{\small(بہار \textenglish{2013}، سوال \textenglish{2})}}
صحافت کی تعریف کریں۔ نیز صحافت کے مختلف نظریات پر بحث کریں۔
\end{sawal}

\begin{hal}
\textbf{تعریف۔} صحافت وہ فن اور پیشہ ہے جس میں حالاتِ حاضرہ کی خبریں جمع کی
جاتی ہیں، ان کی تصدیق و ترتیب کی جاتی ہے، اور انہیں کسی ذریعۂ ابلاغ (اخبار،
رسالہ، ریڈیو، ٹی وی، ویب سائٹ) کے ذریعے عوام تک بروقت پہنچایا جاتا ہے۔ لفظ
''صحافت'' عربی \emph{صحیفہ} سے ہے، یعنی لکھا ہوا ورق۔

\medskip
\textbf{صحافت کے نظریات۔} سیبرٹ، پیٹرسن اور شرام کی مشہور کتاب
\textenglish{\emph{Four Theories of the Press}} (\textenglish{1956}) کے مطابق:

\begin{center}
\begin{tabular}{@{}p{3.4cm} p{5.4cm} p{5.6cm}@{}}
\toprule
\textbf{نظریہ} & \textbf{بنیادی خیال} & \textbf{نتیجہ} \\
\midrule
آمرانہ \newline \textenglish{(Authoritarian)}
  & صحافت ریاست کے تابع ہے & سنسر شپ، اجازت نامہ، سرکاری کنٹرول \\
آزاد خیال \newline \textenglish{(Libertarian)}
  & انسان عقل رکھتا ہے، سچ خود سامنے آ جائے گا & مکمل آزادی، ریاست کی کم سے کم مداخلت \\
سماجی ذمہ داری \newline \textenglish{(Social Responsibility)}
  & آزادی کے ساتھ ذمہ داری بھی & ضابطۂ اخلاق، خود احتسابی، پریس کونسل \\
اشتراکی \newline \textenglish{(Soviet Communist)}
  & ذرائع ابلاغ ریاست و پارٹی کی ملکیت & نظریے کی تبلیغ، نجی ملکیت ممنوع \\
\bottomrule
\end{tabular}
\end{center}

بعد میں دو اور نظریات شامل کیے گئے: \textbf{ترقیاتی نظریہ}
(\textenglish{Development}) --- ترقی پذیر ممالک میں صحافت قومی تعمیر میں
حصہ لے؛ اور \textbf{جمہوری شرکتی نظریہ}
(\textenglish{Democratic Participant}) --- مقامی اور چھوٹے ذرائع ابلاغ کو
عوام کی آواز بننا چاہیے۔

\medskip
\textbf{اسلامی نظریۂ صحافت۔} صداقت، امانت، تحقیق (\emph{تبیّنوا})، غیبت و
بہتان سے اجتناب اور امر بالمعروف --- یعنی آزادی اور اخلاقی حدود دونوں۔

\jawab{صحافت خبروں کے حصول، تصدیق اور ترسیل کا فن ہے۔ اہم نظریات: آمرانہ،
آزاد خیال، سماجی ذمہ داری، اشتراکی، اور بعد ازاں ترقیاتی و جمہوری شرکتی۔
پاکستان میں عملاً \emph{سماجی ذمہ داری} کا نظریہ رائج ہے۔}
\end{hal}

\begin{sawal}{سوال \textenglish{2} \ \textnormal{\small(بہار \textenglish{2013}، سوال \textenglish{1} اور خزاں \textenglish{2013}، سوال \textenglish{1})}}
برصغیر میں صحافت کی ابتداء و ارتقاء کی وضاحت کریں، اور پاکستانی صحافت کے
تاریخی ارتقاء پر نوٹ لکھیں۔
\end{sawal}

\begin{hal}
یہ سوال دونوں پرچوں میں آیا --- سب سے اہم سوال یہی ہے۔ جواب کو تین ادوار میں
بانٹیے۔

\medskip
\textbf{(۱) آغاز (\textenglish{1780--1857})}
\begin{itemize}\setlength{\itemsep}{1pt}
  \item برصغیر کا پہلا اخبار \textbf{\textenglish{Bengal Gazette}}
        (\textenglish{1780}، کلکتہ)، مدیر \textenglish{James Augustus Hicky}۔
  \item پہلا اردو اخبار \textbf{جامِ جہاں نما} (\textenglish{1822}، کلکتہ)۔
  \item \textbf{دہلی اردو اخبار} (\textenglish{1837}) --- مولوی محمد باقر، جو
        \textenglish{1857} میں شہید ہوئے؛ اردو صحافت کے پہلے شہید صحافی۔
\end{itemize}

\textbf{(۲) تحریکِ آزادی کا دور (\textenglish{1857--1947})}
\begin{itemize}\setlength{\itemsep}{1pt}
  \item سر سید احمد خان: \textbf{سائنٹفک سوسائٹی گزٹ} (\textenglish{1866}) اور
        \textbf{تہذیب الاخلاق} (\textenglish{1870})۔
  \item مولانا ظفر علی خان: \textbf{زمیندار} --- مسلمانوں کی توانا آواز۔
  \item مولانا محمد علی جوہر: \textbf{کامریڈ} (\textenglish{1911}، انگریزی) اور
        \textbf{ہمدرد} (\textenglish{1913}، اردو)۔
  \item مولانا ابوالکلام آزاد: \textbf{الہلال} (\textenglish{1912}) اور
        \textbf{البلاغ} (\textenglish{1915})۔
  \item \textbf{انقلاب} (\textenglish{1927}، لاہور)، \textbf{ڈان}
        (\textenglish{1941}، قائدِ اعظم کی رہنمائی میں)، \textbf{نوائے وقت}
        (\textenglish{1940})، \textbf{جنگ} (\textenglish{1939})۔
\end{itemize}

\textbf{(۳) قیامِ پاکستان کے بعد (\textenglish{1947} تا حال)}
\begin{itemize}\setlength{\itemsep}{1pt}
  \item ہجرت کے بعد اخبارات کی پاکستان منتقلی؛ کاغذ، سرمائے اور عملے کی قلت۔
  \item مارشل لا کے ادوار میں سنسر شپ اور \textenglish{PPO 1963} جیسے سخت
        قوانین؛ \textenglish{1988} میں اس کا خاتمہ۔
  \item \textenglish{2002} میں \textenglish{PEMRA} کا قیام --- نجی ٹی وی چینلوں
        کا آغاز اور الیکٹرانک میڈیا کا دھماکہ خیز پھیلاؤ۔
  \item حالیہ دور: ڈیجیٹل صحافت، سوشل میڈیا اور آن لائن اخبارات۔
\end{itemize}

\jawab{برصغیر میں صحافت کا آغاز \textenglish{Bengal Gazette}
(\textenglish{1780}) سے ہوا اور پہلا اردو اخبار جامِ جہاں نما
(\textenglish{1822}) تھا۔ تحریکِ آزادی میں زمیندار، الہلال، کامریڈ اور ہمدرد
نے مرکزی کردار ادا کیا۔ قیامِ پاکستان کے بعد صحافت نے قلت، سنسر شپ اور پھر
\textenglish{2002} کے بعد الیکٹرانک و ڈیجیٹل انقلاب کے مراحل دیکھے۔}
\end{hal}

\begin{sawal}{سوال \textenglish{3} \ \textnormal{\small(خزاں \textenglish{2013}، سوال \textenglish{2})}}
پاکستان کی تخلیق میں سر سید احمد خان اور مولانا ظفر علی خان کی صحافتی خدمات پر
بحث کیجیے۔
\end{sawal}

\begin{hal}
\textbf{سر سید احمد خان (\textenglish{1817--1898})}
\begin{itemize}\setlength{\itemsep}{1pt}
  \item \textbf{سائنٹفک سوسائٹی گزٹ} (\textenglish{1866}) --- دو لسانی؛ مغربی
        علوم کا تعارف اور مسلمانوں کو جدید تعلیم کی ترغیب۔
  \item \textbf{تہذیب الاخلاق} (\textenglish{1870}) --- مسلمانوں کی سماجی و
        فکری اصلاح؛ اردو نثر کو سادہ، مدلل اور مضمون نگاری کے قابل بنایا۔
  \item \textbf{اثر:} \emph{دو قومی نظریے} کی فکری بنیاد رکھی، مسلمانوں کو
        انگریزوں سے کشیدگی کے بجائے تعلیم کی طرف موڑا، اور علی گڑھ تحریک کے
        ذریعے وہ طبقہ پیدا کیا جو بعد میں تحریکِ پاکستان کی قیادت بنا۔
\end{itemize}

\textbf{مولانا ظفر علی خان (\textenglish{1873--1956})}
\begin{itemize}\setlength{\itemsep}{1pt}
  \item \textbf{زمیندار} (لاہور) کے مدیر --- اسے برصغیر کے مؤثر ترین اردو
        اخبار میں بدل دیا۔
  \item بے باک اداریے اور شاعری؛ تحریکِ خلافت اور مسلمانوں کے حقوق کی
        پرزور وکالت۔
  \item اخبار کئی بار بند ہوا، ضمانتیں ضبط ہوئیں اور وہ خود قید ہوئے --- مگر
        مؤقف نہیں بدلا۔
  \item \textbf{اثر:} مسلم لیگ اور قائدِ اعظم کے مؤقف کو عوامی سطح تک پہنچایا؛
        انہیں \emph{بابائے صحافت} کہا جاتا ہے۔
\end{itemize}

\medskip
\textbf{فرق۔} سر سید کی صحافت \emph{اصلاحی و تعلیمی} تھی --- قوم کو تیار کرنا۔
ظفر علی خان کی صحافت \emph{سیاسی و جذباتی} تھی --- قوم کو متحرک کرنا۔ دونوں
مل کر تحریکِ پاکستان کی فکری اور جذباتی بنیاد بنے۔

\jawab{سر سید نے سائنٹفک سوسائٹی گزٹ اور تہذیب الاخلاق کے ذریعے مسلمانوں کی
تعلیمی و فکری اصلاح کی اور دو قومی نظریے کی بنیاد رکھی؛ مولانا ظفر علی خان نے
زمیندار کے ذریعے سیاسی بیداری پیدا کی اور مسلم لیگ کے مؤقف کو عوام تک
پہنچایا۔}
\end{hal}

%---------------------------------------------------------------------
\section*{صحافتی اصناف}
\markboth{صحافتی اصناف}{}

\begin{sawal}{سوال \textenglish{4} \ \textnormal{\small(خزاں \textenglish{2013}، سوال \textenglish{4})}}
اداریے سے کیا مراد ہے؟ نیز اداریے کی مختلف اقسام مثالوں کے ساتھ بیان کریں۔
\end{sawal}

\begin{hal}
\textbf{تعریف۔} اداریہ (\textenglish{Editorial}) وہ تحریر ہے جس میں اخبار کسی
اہم واقعے یا مسئلے پر اپنی \emph{ادارتی رائے} پیش کرتا ہے۔ یہ کسی ایک شخص کی
نہیں بلکہ پورے ادارے کی رائے سمجھی جاتی ہے، اسی لیے عام طور پر اس پر لکھنے
والے کا نام نہیں ہوتا۔ اسے اخبار کی \emph{آواز} یا \emph{ضمیر} کہا جاتا ہے۔

\medskip
\textbf{اجزا:} عنوان، تمہید (واقعے کا بیان)، بحث و تجزیہ، اور نتیجہ یا تجویز۔

\medskip
\textbf{اقسام:}
\begin{center}
\begin{tabular}{@{}p{3.4cm} p{5.4cm} p{5.6cm}@{}}
\toprule
\textbf{قسم} & \textbf{مقصد} & \textbf{مثال} \\
\midrule
تشریحی & واقعے کا پس منظر اور مطلب سمجھانا & بجٹ کی تفصیل کی وضاحت \\
تنقیدی & کسی پالیسی یا اقدام پر گرفت & بجلی کی قیمتوں میں اضافے پر تنقید \\
معلوماتی & حقائق و اعداد و شمار کی فراہمی & مردم شماری کے نتائج \\
ستائشی & کسی اچھے کام کی تحسین & کھیلوں میں قومی کامیابی \\
مہماتی & کسی مقصد کے لیے رائے عامہ ہموار کرنا & شجرکاری یا پولیو مہم \\
تعزیتی & کسی اہم شخصیت کی وفات پر & قومی رہنما کا انتقال \\
\bottomrule
\end{tabular}
\end{center}

\jawab{اداریہ اخبار کی ادارتی رائے ہے جو ادارے کی نمائندگی کرتی ہے اور عموماً
بے نام ہوتی ہے۔ اقسام: تشریحی، تنقیدی، معلوماتی، ستائشی، مہماتی اور تعزیتی۔}
\end{hal}

\begin{sawal}{سوال \textenglish{5} \ \textnormal{\small(بہار \textenglish{2013}، سوال \textenglish{3})}}
فیچر اور اداریے میں فرق واضح کریں۔
\end{sawal}

\begin{hal}
\begin{center}
\begin{tabular}{@{}p{2.8cm} p{5.8cm} p{5.8cm}@{}}
\toprule
 & \textbf{اداریہ} & \textbf{فیچر} \\
\midrule
رائے کس کی & پورے ادارے کی & لکھنے والے کی \\
نام & عموماً بے نام & لکھنے والے کے نام سے \\
مقصد & رائے دینا اور رہنمائی & معلومات، دلچسپی، انسانی پہلو \\
موضوع & حالاتِ حاضرہ کا اہم مسئلہ & کوئی بھی موضوع --- شخصیت، مقام، مسئلہ \\
وقت کی پابندی & لازمی --- تازہ واقعے سے جڑا & کم --- بعد میں بھی پڑھا جا سکتا ہے \\
اسلوب & سنجیدہ، مدلل، مختصر & ادبی، رواں، تفصیلی \\
مقام & ادارتی صفحہ & کوئی بھی صفحہ، تصاویر کے ساتھ \\
طوالت & محدود (چند سو الفاظ) & نسبتاً طویل \\
\bottomrule
\end{tabular}
\end{center}

\medskip
\textbf{مشترک بات:} دونوں سادہ خبر نہیں ہیں --- خبر صرف \emph{کیا ہوا} بتاتی
ہے، جبکہ اداریہ اور فیچر دونوں \emph{کیوں} اور \emph{کیسے} کی طرف جاتے ہیں۔

\jawab{اداریہ ادارے کی بے نام رائے ہے جو تازہ اہم مسئلے پر سنجیدہ اور مختصر
انداز میں لکھی جاتی ہے؛ فیچر لکھنے والے کے نام سے شائع ہونے والی تفصیلی، ادبی
اور دلچسپ تحریر ہے جو وقت کی پابند نہیں۔}
\end{hal}

\begin{sawal}{سوال \textenglish{6} \ \textnormal{\small(بہار \textenglish{2013}، سوال \textenglish{5})}}
کالم کی اقسام بیان کریں۔ نیز موجودہ دور میں کالم نویسی کی ضرورت کی وضاحت کریں۔
\end{sawal}

\begin{hal}
\textbf{تعریف۔} کالم اخبار میں باقاعدگی سے شائع ہونے والی وہ تحریر ہے جو کسی
مخصوص لکھنے والے کے نام اور عموماً تصویر کے ساتھ ایک مقررہ جگہ پر آتی ہے، اور
جس میں وہ اپنی ذاتی رائے اور اندازِ نظر پیش کرتا ہے۔

\medskip
\textbf{اقسام:}
\begin{center}
\begin{tabular}{@{}p{3.4cm} p{11.0cm}@{}}
\toprule
سیاسی کالم & حالاتِ حاضرہ اور سیاست پر تجزیہ --- سب سے مقبول قسم \\
طنز و مزاح & ہلکے پھلکے انداز میں سنجیدہ بات؛ مثلاً ابنِ انشا کا اسلوب \\
ادبی کالم & کتابوں، شاعری اور ادبی سرگرمیوں پر \\
سماجی کالم & معاشرتی مسائل --- تعلیم، صحت، خواتین، نوجوان \\
کھیلوں کا کالم & کھیل اور کھلاڑیوں پر تبصرہ \\
مذہبی کالم & دینی موضوعات اور اصلاحِ معاشرہ \\
\bottomrule
\end{tabular}
\end{center}

\medskip
\textbf{موجودہ دور میں ضرورت:}
\begin{itemize}\setlength{\itemsep}{1pt}
  \item خبروں کے سیلاب میں قاری کو \emph{تجزیہ} چاہیے --- کالم یہی دیتا ہے۔
  \item عوامی مسائل کو ایوانِ اقتدار تک پہنچانے کا مؤثر ذریعہ۔
  \item رائے عامہ کی تشکیل اور جمہوری بحث کو زندہ رکھنا۔
  \item قاری کا اخبار سے ذاتی تعلق --- بہت سے لوگ اخبار کسی مخصوص کالم نگار
        کے لیے خریدتے ہیں۔
  \item اردو زبان و ادب کی ترویج۔
\end{itemize}

\jawab{کالم مخصوص لکھاری کی باقاعدہ، نام کے ساتھ شائع ہونے والی رائے پر مبنی
تحریر ہے۔ اقسام: سیاسی، طنز و مزاح، ادبی، سماجی، کھیل اور مذہبی۔ آج اس کی
ضرورت خبروں کے انبار میں تجزیہ فراہم کرنے، عوامی مسائل اجاگر کرنے اور رائے
عامہ بنانے کے لیے ہے۔}
\end{hal}

\begin{sawal}{سوال \textenglish{7} \ \textnormal{\small(بہار \textenglish{2013}، سوال \textenglish{7} اور خزاں \textenglish{2013}، سوال \textenglish{3})}}
سب ایڈیٹر کے اوصاف اور فرائض پر بحث کیجیے۔
\end{sawal}

\begin{hal}
یہ سوال بھی دونوں پرچوں میں آیا۔ \textbf{سب ایڈیٹر} وہ شخص ہے جو رپورٹر کی
بھیجی ہوئی خام خبر کو اشاعت کے قابل بناتا ہے۔ اسے اخبار کا \emph{دربان}
کہا جاتا ہے --- جو کچھ چھپتا ہے، اسی کی چھلنی سے گزر کر چھپتا ہے۔

\medskip
\textbf{فرائض:}
\begin{itemize}\setlength{\itemsep}{1pt}
  \item خبر کی زبان، املا اور قواعد درست کرنا۔
  \item حقائق، ناموں، تاریخوں اور اعداد و شمار کی جانچ۔
  \item مناسب اور کشش رکھنے والی \textbf{سرخی} لگانا --- یہ سب سے اہم فرض ہے۔
  \item غیر ضروری، دہرایا ہوا یا قابلِ اعتراض مواد نکالنا۔
  \item خبر کو مقررہ جگہ کے مطابق مختصر کرنا۔
  \item خبر کی اہمیت کے مطابق صفحے پر مقام کا تعین۔
  \item قانونی پہلو دیکھنا --- ہتکِ عزت یا توہین سے بچاؤ۔
  \item ڈیڈ لائن کی پابندی۔
\end{itemize}

\textbf{اوصاف:}
\begin{itemize}\setlength{\itemsep}{1pt}
  \item زبان پر مکمل عبور (اردو اور انگریزی دونوں)۔
  \item وسیع عمومی معلومات اور حالاتِ حاضرہ سے آگاہی۔
  \item تیز رفتاری اور دباؤ میں کام کرنے کی صلاحیت۔
  \item غیر جانبداری اور دیانت۔
  \item صحافتی قوانین اور ضابطۂ اخلاق سے واقفیت۔
  \item باریک بینی --- ایک غلط ہندسہ پورے اخبار کی ساکھ لے ڈوبتا ہے۔
\end{itemize}

\jawab{سب ایڈیٹر خام خبر کو اشاعت کے قابل بناتا ہے: زبان درست کرنا، حقائق
جانچنا، سرخی لگانا، مواد مختصر کرنا اور قانونی پہلو دیکھنا۔ اوصاف میں زبان پر
عبور، وسیع معلومات، رفتار، غیر جانبداری اور باریک بینی شامل ہیں۔}
\end{hal}

\begin{ghalti}
سب ایڈیٹر اور رپورٹر کے فرائض خلط ملط نہ کیجیے۔ \textbf{رپورٹر} خبر
\emph{لاتا} ہے (میدان میں جا کر)، \textbf{سب ایڈیٹر} خبر کو \emph{سنوارتا} ہے
(دفتر میں بیٹھ کر)۔ سرخی لگانا ہمیشہ سب ایڈیٹر کا کام ہے۔
\end{ghalti}

%---------------------------------------------------------------------
\section*{قوانین، اخلاقیات اور آزادی}
\markboth{قوانین، اخلاقیات اور آزادی}{}

\begin{sawal}{سوال \textenglish{8} \ \textnormal{\small(بہار \textenglish{2013}، سوال \textenglish{4} اور خزاں \textenglish{2013}، سوال \textenglish{5})}}
صحافت میں ضابطۂ اخلاق کی ضرورت اور پاکستانی صحافت میں ضابطۂ اخلاق پر تبصرہ
کیجیے۔
\end{sawal}

\begin{hal}
\textbf{ضابطۂ اخلاق کیا ہے؟} وہ غیر تحریری یا تحریری اصول جن کا پابند رہنا
صحافی اپنے پیشے کے تقاضے کے طور پر قبول کرتا ہے۔ قانون ریاست نافذ کرتی ہے،
ضابطۂ اخلاق صحافی خود پر لاگو کرتا ہے --- یہی بنیادی فرق ہے۔

\medskip
\textbf{بنیادی اصول:}
\begin{itemize}\setlength{\itemsep}{1pt}
  \item \textbf{صداقت اور تصدیق} --- خبر شائع کرنے سے پہلے تحقیق۔
  \item \textbf{غیر جانبداری} --- خبر اور رائے کو الگ رکھنا۔
  \item \textbf{توازن} --- دونوں فریقوں کا مؤقف۔
  \item \textbf{نجی زندگی کا احترام} اور متاثرین کی عزتِ نفس۔
  \item \textbf{ذرائع کا تحفظ} --- خبر دینے والے کا نام ظاہر نہ کرنا۔
  \item \textbf{غلطی کی تصحیح} --- غلط خبر پر اعتراف اور تردید کی اشاعت۔
  \item \textbf{مفادات کا ٹکراؤ} اور لفافہ صحافت سے اجتناب۔
  \item فرقہ وارانہ، نسلی اور صنفی منافرت سے گریز۔
\end{itemize}

\textbf{ضرورت کیوں؟} کیونکہ صحافت کا اثر بہت وسیع ہے: ایک غیر مصدقہ خبر کسی
کی زندگی، ساکھ یا امنِ عامہ کو تباہ کر سکتی ہے۔ نیز اگر صحافت خود کو ضابطے کا
پابند نہ کرے تو ریاست کو سخت قانون سازی کا جواز مل جاتا ہے --- یعنی ضابطۂ
اخلاق دراصل \emph{آزادیٔ صحافت کا محافظ} ہے۔

\medskip
\textbf{پاکستان میں صورتِ حال۔} یہاں ضابطے موجود ہیں ---
\textenglish{APNS}، \textenglish{CPNE}، \textenglish{PFUJ} کے ضابطے اور
\textbf{پریس کونسل آف پاکستان} کا ''\textenglish{Ethical Code of Practice}''
(\textenglish{2002})، نیز \textenglish{PEMRA} کا الیکٹرانک میڈیا ضابطہ۔ مگر
عملی مسائل نمایاں ہیں: ریٹنگ کی دوڑ، بریکنگ نیوز کا دباؤ، غیر مصدقہ خبریں،
لفافہ صحافت، اور نجی زندگی کی خلاف ورزی۔ ادارے موجود ہیں، مگر ان کا نفاذ کمزور
ہے۔

\jawab{ضابطۂ اخلاق صحافی کا خود پر عائد کردہ اصولی نظام ہے: صداقت، غیر
جانبداری، توازن، نجی زندگی کا احترام اور ذرائع کا تحفظ۔ پاکستان میں
\textenglish{APNS}، \textenglish{CPNE}، \textenglish{PFUJ} اور پریس کونسل کے
ضابطے موجود ہیں، مگر ریٹنگ کے دباؤ اور کمزور نفاذ کے سبب عمل درآمد ناقص ہے۔}
\end{hal}

\begin{sawal}{سوال \textenglish{9} \ \textnormal{\small(بہار \textenglish{2013}، سوال \textenglish{4})}}
صحافتی قوانین پر بحث کیجیے۔
\end{sawal}

\begin{hal}
\textbf{آئینی بنیاد۔} آئینِ پاکستان \textenglish{1973} کا \textbf{آرٹیکل
\textenglish{19}} اظہارِ رائے اور آزادیٔ صحافت کی ضمانت دیتا ہے --- مگر
''معقول پابندیوں'' کے ساتھ: اسلام کی عظمت، ملکی سلامتی، امنِ عامہ، اخلاق اور
عدالت کی توہین سے متعلق۔ نیز \textbf{آرٹیکل \textenglish{19-A}}
(\textenglish{2010}) معلومات تک رسائی کا حق دیتا ہے۔

\medskip
\begin{center}
\begin{tabular}{@{}p{6.4cm} p{8.0cm}@{}}
\toprule
\textbf{قانون} & \textbf{موضوع} \\
\midrule
\textenglish{Press and Publications Ordinance 1963} & سخت ترین دور؛ \textenglish{1988} میں منسوخ \\
\textenglish{RPPPO 2002} & پریس اور اشاعت کی رجسٹریشن \\
\textenglish{PEMRA Ordinance 2002} & نجی ریڈیو و ٹی وی چینلوں کا ضابطہ \\
\textenglish{Press Council of Pakistan Ordinance 2002} & ضابطۂ اخلاق اور شکایات \\
\textenglish{Defamation Ordinance 2002} & ہتکِ عزت \\
\textenglish{Freedom of Information Ordinance 2002} & سرکاری معلومات تک رسائی \\
\textenglish{PECA 2016} & برقی جرائم اور آن لائن مواد \\
\textenglish{Right of Access to Information Act 2017} & معلومات کا حق \\
\bottomrule
\end{tabular}
\end{center}

\medskip
\textbf{بنیادی کشمکش۔} ہر معاشرے کو دو چیزوں میں توازن رکھنا پڑتا ہے: عوام کا
\emph{جاننے کا حق}، اور فرد کی \emph{عزت و نجی زندگی} نیز ریاست کی سلامتی۔
صحافتی قوانین اسی توازن کی کوشش ہیں --- اور تنقید یہ ہوتی ہے کہ بعض اوقات
''معقول پابندی'' کی تعریف بہت وسیع کر لی جاتی ہے۔

\jawab{بنیاد آرٹیکل \textenglish{19} اور \textenglish{19-A} ہیں۔ اہم قوانین:
\textenglish{PPO 1963} (منسوخ)، \textenglish{RPPPO 2002}،
\textenglish{PEMRA 2002}، پریس کونسل آرڈیننس \textenglish{2002}، ہتکِ عزت
آرڈیننس \textenglish{2002}، \textenglish{PECA 2016} اور معلومات تک رسائی کا
ایکٹ \textenglish{2017}۔}
\end{hal}

\begin{sawal}{سوال \textenglish{10} \ \textnormal{\small(بہار \textenglish{2013}، سوال \textenglish{6})}}
کیا آپ سمجھتے ہیں کہ پاکستانی میڈیا آزاد ہے؟ مثالوں سے وضاحت کریں۔
\end{sawal}

\begin{hal}
یہ \emph{رائے} کا سوال ہے۔ ممتحن متوازن جواب چاہتا ہے --- دونوں پہلو، پھر
اپنی رائے۔

\medskip
\textbf{آزادی کے دلائل:}
\begin{itemize}\setlength{\itemsep}{1pt}
  \item \textenglish{2002} کے بعد درجنوں نجی چینل؛ ریاستی اجارہ داری ختم۔
  \item ٹاک شوز میں حکومتی پالیسیوں پر کھلی تنقید۔
  \item تحقیقاتی صحافت نے بڑے مالی و انتظامی معاملات بے نقاب کیے۔
  \item عدلیہ کے کئی فیصلے میڈیا کے حق میں آئے؛ آرٹیکل \textenglish{19}
        کا تحفظ۔
  \item سوشل میڈیا نے روایتی دربانی (\textenglish{gatekeeping}) کمزور کر دی۔
\end{itemize}

\textbf{پابندیوں کے دلائل:}
\begin{itemize}\setlength{\itemsep}{1pt}
  \item بعض موضوعات پر عملی خود ساختہ سنسر شپ۔
  \item سرکاری اشتہارات کا معاشی دباؤ کے طور پر استعمال۔
  \item مالکان کے کاروباری مفادات کا ادارتی پالیسی پر اثر۔
  \item صحافیوں کو دھمکیاں اور تحفظ کا فقدان۔
  \item ریٹنگ کی دوڑ --- یہ آزادی کو اندر سے نقصان پہنچاتی ہے۔
\end{itemize}

\medskip
\textbf{متوازن نتیجہ۔} پاکستانی میڈیا \emph{قانونی طور پر} آزاد ہے اور ماضی کے
مقابلے میں کہیں زیادہ آزاد ہے، مگر \emph{عملی طور پر} معاشی دباؤ، تحفظ کے فقدان
اور تجارتی مفادات کی وجہ سے یہ آزادی مکمل نہیں۔ حقیقی آزادی صرف قانون سے نہیں،
پیشہ ورانہ ذمہ داری اور معاشی خود مختاری سے آتی ہے۔

\jawab{پاکستانی میڈیا آئینی اعتبار سے آزاد اور \textenglish{2002} کے بعد بہت
زیادہ خود مختار ہے، مگر معاشی دباؤ، سرکاری اشتہارات، مالکان کے مفادات اور
صحافیوں کے تحفظ کے فقدان کے باعث یہ آزادی عملاً محدود ہے۔}
\end{hal}

\begin{ghalti}
اس قسم کے ''کیا آپ سمجھتے ہیں'' والے سوال میں صرف ایک طرف کا مؤقف نہ لکھیے۔
پہلے \textbf{دونوں} پہلو دلائل کے ساتھ دیجیے، پھر اپنی رائے۔ یکطرفہ جواب آدھے
نمبر لیتا ہے۔
\end{ghalti}

%---------------------------------------------------------------------
\section*{صحافت، معاشرہ اور فن}
\markboth{صحافت، معاشرہ اور فن}{}

\begin{sawal}{سوال \textenglish{11} \ \textnormal{\small(خزاں \textenglish{2013}، سوال \textenglish{7})}}
معاشرے کی ترقی میں صحافت کی اہمیت اجاگر کریں، اور جمہوریت کی بقا کے لیے صحافت
کی ضرورت پر لکھیں۔
\end{sawal}

\begin{hal}
\textbf{معاشرے میں صحافت کے کام:}
\begin{itemize}\setlength{\itemsep}{1pt}
  \item \textbf{اطلاع} --- عوام کو حالات سے باخبر رکھنا۔
  \item \textbf{تعلیم و شعور} --- صحت، ماحول، حقوق اور فرائض کی آگاہی۔
  \item \textbf{نگرانی} --- حکومت اور اداروں کے احتساب کا ذریعہ۔
  \item \textbf{رائے عامہ کی تشکیل} اور قومی مکالمے کا پلیٹ فارم۔
  \item \textbf{ثقافتی ورثے کی منتقلی} اور زبان کی ترویج۔
  \item \textbf{تفریح} اور اشتہار کے ذریعے معاشی سرگرمی۔
\end{itemize}

\medskip
\textbf{جمہوریت اور صحافت --- چار بنیادی روابط:}
\begin{enumerate}\setlength{\itemsep}{1pt}
  \item \textbf{باخبر ووٹر۔} جمہوریت کا انحصار باشعور رائے دہندہ پر ہے، اور
        شعور اطلاع سے آتا ہے۔ بے خبر ووٹ جمہوریت نہیں، ہجوم ہے۔
  \item \textbf{چوتھا ستون۔} مقننہ، عدلیہ اور انتظامیہ کے بعد صحافت وہ ادارہ
        ہے جو باقی تینوں پر نظر رکھتا ہے۔
  \item \textbf{شفافیت اور احتساب۔} بدعنوانی کا سب سے بڑا دشمن اشاعت ہے۔
  \item \textbf{اختلافِ رائے کا پُرامن اظہار۔} جہاں بولنے کا راستہ کھلا ہو،
        وہاں تشدد کا راستہ تنگ ہوتا ہے۔
\end{enumerate}

\medskip
\textbf{شرط۔} یہ سب تبھی ممکن ہے جب صحافت خود آزاد، ذمہ دار اور غیر جانبدار
ہو۔ جانبدار یا بکاؤ صحافت جمہوریت کو سہارا دینے کے بجائے نقصان پہنچاتی ہے۔

\jawab{صحافت معاشرے میں اطلاع، تعلیم، نگرانی، رائے عامہ کی تشکیل اور ثقافتی
ترسیل کا کام کرتی ہے۔ جمہوریت کے لیے وہ ناگزیر ہے کیونکہ باخبر ووٹر پیدا کرتی
ہے، چوتھے ستون کے طور پر احتساب کرتی ہے، شفافیت لاتی ہے اور اختلافِ رائے کو
پُرامن راستہ دیتی ہے۔}
\end{hal}

\begin{sawal}{سوال \textenglish{12} \ \textnormal{\small(خزاں \textenglish{2013}، سوال \textenglish{6})}}
موجودہ دور کی صحافت میں جدید رجحانات کے اثرات پر روشنی ڈالیں۔
\end{sawal}

\begin{hal}
\textbf{اہم جدید رجحانات:}
\begin{itemize}\setlength{\itemsep}{1pt}
  \item \textbf{ڈیجیٹل اور آن لائن صحافت} --- اخبار کی ویب سائٹ اور ایپ۔
  \item \textbf{سوشل میڈیا} --- خبر کی ترسیل اور اب خبر کا ماخذ بھی۔
  \item \textbf{شہری صحافت} (\textenglish{Citizen Journalism}) --- ہر
        موبائل رکھنے والا ممکنہ رپورٹر۔
  \item \textbf{\textenglish{24/7} خبری چینل} اور بریکنگ نیوز کلچر۔
  \item \textbf{کثیر الوسائل صحافت} (\textenglish{Multimedia}) --- متن،
        ویڈیو، انفوگرافک ایک ساتھ۔
  \item \textbf{ڈیٹا صحافت} اور \textenglish{AI} کا بڑھتا استعمال۔
\end{itemize}

\medskip
\begin{center}
\begin{tabular}{@{}p{7.0cm} p{7.0cm}@{}}
\toprule
\textbf{مثبت اثرات} & \textbf{منفی اثرات} \\
\midrule
خبر کی رفتار --- لمحوں میں دنیا بھر تک & تصدیق کا وقت نہ ملنا؛ غلط خبر کا پھیلاؤ \\
ہر شخص کو اظہار کا موقع & جعلی خبریں اور افواہیں \\
اشاعت کی لاگت میں زبردست کمی & اخبارات کی معاشی مشکلات اور بندش \\
عالمی رسائی & مقامی صحافت کی کمزوری \\
باہمی رابطہ --- قاری جواب دے سکتا ہے & ریٹنگ اور ''کلک'' کی دوڑ؛ سنسنی خیزی \\
معلومات کے ذخائر تک رسائی & نجی زندگی کی خلاف ورزی \\
\bottomrule
\end{tabular}
\end{center}

\medskip
\textbf{نتیجہ۔} جدید ٹیکنالوجی نے صحافت کی \emph{رفتار اور رسائی} بے پناہ بڑھا
دی ہے، مگر \emph{تصدیق اور ذمہ داری} کو مشکل بنا دیا ہے۔ آج کے صحافی کا اصل
امتحان یہ ہے کہ وہ رفتار کی دوڑ میں صداقت قربان نہ کرے۔

\jawab{جدید رجحانات میں ڈیجیٹل صحافت، سوشل میڈیا، شہری صحافت،
\textenglish{24/7} چینل اور ڈیٹا صحافت شامل ہیں۔ ان سے رفتار، رسائی اور
شراکت بڑھی، مگر جعلی خبریں، سنسنی خیزی، اخبارات کی معاشی مشکلات اور نجی زندگی
کی خلاف ورزی جیسے مسائل پیدا ہوئے۔}
\end{hal}

\begin{sawal}{سوال \textenglish{13} \ \textnormal{\small(بہار \textenglish{2013}، سوال \textenglish{8}--الف)}}
طباعت کی اقسام پر نوٹ لکھیں۔
\end{sawal}

\begin{hal}
\begin{center}
\begin{tabular}{@{}p{3.4cm} p{5.6cm} p{5.4cm}@{}}
\toprule
\textbf{قسم} & \textbf{طریقۂ کار} & \textbf{استعمال} \\
\midrule
لیٹر پریس \newline \textenglish{(Letterpress)}
  & ابھرے ہوئے حروف پر سیاہی، پھر کاغذ پر دباؤ & پرانا طریقہ؛ اب تقریباً متروک \\
آفسیٹ \newline \textenglish{(Offset Litho)}
  & پلیٹ سے ربڑ کے کمبل پر، پھر کاغذ پر & اخبارات و کتب --- آج سب سے عام \\
گریوئیر \newline \textenglish{(Gravure)}
  & گہرے کھدے خانوں میں سیاہی بھر کر & رسائل، پیکنگ؛ بہت بڑی تعداد \\
فلیکسو گرافی \newline \textenglish{(Flexography)}
  & لچکدار ربڑ پلیٹ سے & پیکنگ، تھیلے، لیبل \\
سکرین پرنٹنگ & باریک جالی سے سیاہی گزار کر & کپڑا، بینر، پوسٹر \\
ڈیجیٹل پرنٹنگ & براہِ راست فائل سے، پلیٹ کے بغیر & کم تعداد، فوری کام \\
\bottomrule
\end{tabular}
\end{center}

\medskip
\textbf{ایک اہم نکتہ۔} اخباری صنعت میں لیٹر پریس سے \textbf{آفسیٹ} کی طرف
منتقلی ایک انقلاب تھی: چھپائی تیز، سستی اور بہتر ہوئی، رنگین اشاعت ممکن ہوئی،
اور اخبار کے صفحات و تصاویر کا معیار یکسر بدل گیا۔

\jawab{اہم اقسام: لیٹر پریس، آفسیٹ، گریوئیر، فلیکسو گرافی، سکرین پرنٹنگ اور
ڈیجیٹل پرنٹنگ۔ اخبارات آج بنیادی طور پر آفسیٹ پر چھپتے ہیں۔}
\end{hal}

\begin{sawal}{سوال \textenglish{14} \ \textnormal{\small(بہار \textenglish{2013}، سوال \textenglish{8}--ب)}}
پروپیگنڈہ پر نوٹ لکھیں۔
\end{sawal}

\begin{hal}
\textbf{تعریف۔} پروپیگنڈہ معلومات، دلائل، افواہوں یا علامتوں کا وہ منظم اور
بامقصد استعمال ہے جس کا مقصد لوگوں کی رائے اور رویّے کو کسی مخصوص سمت میں
بدلنا ہو --- خواہ اس کے لیے حقائق کو توڑا مروڑا ہی کیوں نہ جائے۔

\medskip
\textbf{اقسام (ماخذ کے اعتبار سے):}
\begin{itemize}\setlength{\itemsep}{1pt}
  \item \textbf{سفید} --- ماخذ ظاہر اور معلومات درست (سرکاری آگاہی مہم)۔
  \item \textbf{سیاہ} --- ماخذ جھوٹا اور معلومات گمراہ کن (دشمن کے نام سے
        جعلی پیغام)۔
  \item \textbf{سرمئی} --- ماخذ مبہم، صداقت مشکوک۔
\end{itemize}

\textbf{عام تکنیکیں:}
\begin{center}
\begin{tabular}{@{}p{4.4cm} p{9.8cm}@{}}
\toprule
نام دھرنا (\textenglish{Name Calling}) & مخالف پر منفی لقب چسپاں کرنا \\
خوشنما الفاظ (\textenglish{Glittering Generality}) & مبہم مگر پرکشش نعرے \\
منتقلی (\textenglish{Transfer}) & کسی محترم علامت کا سہارا لینا \\
شہادت (\textenglish{Testimonial}) & مشہور شخصیت سے تائید دلوانا \\
عام آدمی (\textenglish{Plain Folks}) & ''میں تو آپ ہی میں سے ہوں'' \\
یکطرفہ پتے (\textenglish{Card Stacking}) & صرف موافق حقائق پیش کرنا \\
ہجوم کا دباؤ (\textenglish{Band Wagon}) & ''سب یہی کر رہے ہیں'' \\
\bottomrule
\end{tabular}
\end{center}

\medskip
\textbf{صحافت اور پروپیگنڈے کا فرق۔} صحافت کا مقصد \emph{اطلاع} ہے اور اس کی
کسوٹی \emph{صداقت}؛ پروپیگنڈے کا مقصد \emph{قائل کرنا} ہے اور اس کی کسوٹی
\emph{اثر}۔ اسی لیے پروپیگنڈہ حقائق کے چناؤ میں یکطرفہ ہوتا ہے، جبکہ صحافت پر
توازن لازم ہے۔

\jawab{پروپیگنڈہ رائے اور رویّے بدلنے کے لیے معلومات کا منظم و بامقصد استعمال
ہے۔ اقسام: سفید، سیاہ، سرمئی۔ تکنیکوں میں نام دھرنا، خوشنما الفاظ، منتقلی،
شہادت، عام آدمی، یکطرفہ پتے اور ہجوم کا دباؤ شامل ہیں۔ صحافت اطلاع دیتی ہے،
پروپیگنڈہ قائل کرتا ہے۔}
\end{hal}
